\newpage{}
\section{Distribution Field and Observer Operators}
\label{sec:distribution-field-and-observer-operators}

Section~\ref{sec:edge-forecast-head} defines an edge-local predictive
distribution
\[
  {}_{\psi}\mathfrak{D}_{u\,v}^{c\,i}
  \in
  \mathcal{P}(\R^{D_f})
\]
for each directed edge $e=(u,v)$, forecast channel $c$, and forecast step $i$.
This section defines the readout layer around those distributions.

An observer is a declared way of reading something: a predictive distribution,
a scalar marginal, a normalized edge field, a node-lifted field, or a residual
field.  Observers do not replace the object they read.  They expose one useful
quantity from it.

\calloutbox{Observer principle.}{\\%
  The predictive distribution is the uncertainty-bearing object.  An observer
  is a readout.  Expectations, variances, quantiles, raw-activity expectations,
  NodeLift inputs, residuals, and control-prepared fields are all different
  views of the same underlying forecasting interface.
}

\subsection{Edge Distribution Field}
\label{subsec:edge-distribution-field}

For a fixed forecast channel $c$ and forecast step $i$, the
\emph{edge distribution field} is the edge-indexed collection
\[
  {}_{\psi}\mathfrak{D}_{\mathcal{E}}^{c\,i}
  =
  \left\{
    {}_{\psi}\mathfrak{D}_{u\,v}^{c\,i}
    \in
    \mathcal{P}(\R^{D_f})
    :
    (u,v)\in\mathcal{E}
  \right\}
  \in
  \left(\mathcal{P}(\R^{D_f})\right)^L.
\]
If the model is evaluated repeatedly over time, each local distribution may be
written
\[
  {}_{\psi}\mathfrak{D}_{u\,v}^{c\,i}(t),
  \qquad
  t\in\mathcal{T},
\]
and the distribution field itself becomes time-indexed:
\[
  t
  \longmapsto
  {}_{\psi}\mathfrak{D}_{\mathcal{E}}^{c\,i}(t).
\]

The field stores uncertainty.  It is not yet a point forecast, node state,
trading score, or action.  Those objects appear only after an observer reads the
field.

\subsection{Scalar Marginals}
\label{subsec:scalar-marginals}

Observers may read the full future-target distribution or a selected scalar
marginal.  For any exposed target coordinate $d$, define the coordinate
projection
\[
  \pi_d:\R^{D_f}\rightarrow\R,
  \qquad
  \pi_d(\vect{f}) = f^d.
\]
The scalar marginal distribution for coordinate $d$ is the pushforward
\[
  {}_{\psi}\mathfrak{D}_{(d),u\,v}^{c\,i}
  =
  (\pi_d)_{\#}
  {}_{\psi}\mathfrak{D}_{u\,v}^{c\,i}
  \in
  \mathcal{P}(\R).
\]

The scalar edge-return marginal is the special case where $d$ is the selected
return coordinate.  Using the semantic label $(r)$,
\[
  {}_{\psi}\mathfrak{D}_{(r),u\,v}^{c\,i}
  =
  (\pi_{(r)})_{\#}
  {}_{\psi}\mathfrak{D}_{u\,v}^{c\,i}
  \in
  \mathcal{P}(\R).
\]
Thus the scalar return marginal does not replace the full predictive
distribution.  It is one observer input derived from it.

\subsection{Observer Operators}
\label{subsec:observer-operators}

For the chosen observer input dimension,
\[
  d_{\mathcal{O}}\in\{1,D_f\},
  \qquad
  \mathcal{O}
  :
  \mathcal{P}(\R^{d_{\mathcal{O}}})
  \rightarrow
  \mathcal{Y},
\]
where $d_{\mathcal{O}}=1$ denotes a scalar observer and
$d_{\mathcal{O}}=D_f$ denotes an observer on the full future target vector.
The output space $\mathcal{Y}$ depends on the observer.

For
\[
  \mathfrak{D}\in\mathcal{P}(\R^{D_f}),
  \qquad
  \vect{F}\sim\mathfrak{D},
\]
useful coordinatewise observers include
\begin{align*}
  \mathcal{O}_{\E}[\mathfrak{D}]
    &=
    \E_{\mathfrak{D}}[\vect{F}]
    &&\in\R^{D_f}
    &&\text{(expected value)},\\
  \mathcal{O}_{\Var}[\mathfrak{D}]
    &=
    \Var_{\mathfrak{D}}[\vect{F}]
    &&\in\R^{D_f}
    &&\text{(coordinatewise variance)},\\
  \mathcal{O}_{\Quant\,\tau}[\mathfrak{D}]
    &=
    \Quant_{\tau}(\mathfrak{D})
    &&\in\R^{D_f}
    &&\text{(coordinatewise $\tau$-quantile)},\\
  \mathcal{O}_{\ES\,\tau}[\mathfrak{D}]
    &=
    \ES_{\tau}(\mathfrak{D})
    &&\in\R^{D_f}
    &&\text{(coordinatewise expected shortfall)},\\
  \mathcal{O}_{\Risk}[\mathfrak{D}]
    &\in \R^m
    &&
    &&\text{(generic risk observer)}.
\end{align*}
Applied to a scalar marginal
\[
  {}_{\psi}\mathfrak{D}_{(d),u\,v}^{c\,i}
  \in
  \mathcal{P}(\R),
\]
the same named observer families return scalar outputs unless a larger output
space is explicitly declared.

For the expected-shortfall observer we follow the coherent-risk treatment of
Acerbi and Tasche~\cite{acerbi2002coherence}.  For mixture distributions,
coordinatewise quantiles and expected shortfall are computed numerically from
the scalar marginal distribution, unless a closed-form special case is
explicitly available.

In plain terms, observers do not create the distribution field.  They read the
field.  The expected value attached to an instrument edge is therefore one
observer among several, not the whole predictive object.

\subsection{Distribution Calibration Diagnostics}
\label{subsec:distribution-calibration-diagnostics}

The MDN loss trains a predictive density, but low training loss alone does not
guarantee calibrated uncertainty.  During evaluation, the distribution field
should be checked with held-out negative log-likelihood, finite-value checks,
scale statistics, mixture-weight entropy, component usage, coverage of
predictive intervals, probability integral transform diagnostics for scalar
marginals, and error statistics of observer readouts such as expectations,
quantiles, and expected shortfall.

These diagnostics keep validation tied to the full predictive distribution
rather than only to a point-estimate error such as mean squared error.

\subsection{Scalar Edge-Return Field Observers}
\label{subsec:scalar-edge-return-field-observers}

Applying a scalar observer pointwise over the edge distribution field gives a
scalar edge field.  For edge $e=(u,v)$, the scalar edge-return marginal gives
the mean-return readout
\[
  \widehat{r}_{e}^{c\,i}
  :=
  \mathcal{O}_{\E}
  \left[
    {}_{\psi}\mathfrak{D}_{(r),u\,v}^{c\,i}
  \right]
  \in\R.
\]
A simple variance-based confidence rule is
\[
  \gamma_{e}^{c\,i}
  :=
  \frac{1}{
    \mathcal{O}_{\Var}
    \left[
      {}_{\psi}\mathfrak{D}_{(r),u\,v}^{c\,i}
    \right]
    +
    \delta
  }
  \in\Rpos,
  \qquad
  \delta\in\Rpos.
\]
High predictive variance lowers the confidence weight; the floor $\delta$
prevents a zero denominator.

For a fixed channel and forecast step, one may suppress $(c,i)$ and write
\[
  \widehat{r}_e
  :=
  \widehat{r}_{e}^{c\,i},
  \qquad
  \gamma_e
  :=
  \gamma_{e}^{c\,i}.
\]
This is the notation used when a scalar price-coordinate field is passed to
NodeLift or to a price-residual diagnostic.

The same pattern gives other scalar edge fields:
\[
  \widehat{r}_{e,\tau}^{c\,i}
  :=
  \mathcal{O}_{\Quant\,\tau}
  \left[
    {}_{\psi}\mathfrak{D}_{(r),u\,v}^{c\,i}
  \right],
\]
\[
  \widehat{r}_{e,\mathrm{ES},\tau}^{c\,i}
  :=
  \mathcal{O}_{\ES\,\tau}
  \left[
    {}_{\psi}\mathfrak{D}_{(r),u\,v}^{c\,i}
  \right].
\]
Thus the scalar edge field need not always be a mean field.  It may be a
quantile field, downside-risk field, or another policy-specific readout.

\subsection{Input Field Observers for NodeLift}
\label{subsec:edge-feature-field-observers-for-nodelift}

NodeLift in Section~\ref{sec:edges-to-nodes} consumes observed input
edge-coordinate fields.

For an observed input edge window, the raw normalized coordinate observer is
\[
  \mathcal{O}_{x,d}
  \left[
    \tensor{x}_{u\,v}^{c\,i}(t)
  \right]
  =
  x_{u\,v}^{c\,i\,d}(t)
  \in\R.
\]
Collecting over usable edges gives the input edge-coordinate field
\[
  \vect{x}_{\mathcal{E}}^{c\,i\,d}(t)
  =
  \left(
    x_e^{c\,i\,d}(t)
  \right)_{e\in\mathcal{E}_{\mathrm{use}}}
  \in
  \R^{L_{\mathrm{use}}}.
\]
This is the object consumed by feature-wise NodeLift.  Predictive and
supervised readouts can still be defined as observers, but if a later extension
lifts them it should name the source field explicitly.

For price coordinates, this normalized coordinate field can be lifted directly
through the signed incidence relation.  For activity coordinates, one must read
or reconstruct raw nonnegative activity before endpoint aggregation.

\subsection{Raw-Activity Observers}
\label{subsec:raw-activity-observers}

Activity coordinates are stored in normalized form as
\[
  X=\log(1+A),
  \qquad
  A\in\Rnonneg.
\]
For observed normalized input data, the raw-activity observer is direct:
\[
  \mathcal{O}_{\mathrm{rawact}}
  \left[
    x_e^{c\,i\,d}
  \right]
  =
  \exp
  \left(
    x_e^{c\,i\,d}
  \right)
  -
  1
  \in\Rnonneg,
  \qquad
  d\in\mathcal{J}_{\mathrm{act}}.
\]

For predictive distributions, the expectation should be taken after the inverse
activity transform.  If
\[
  X
  \sim
  {}_{\psi}\mathfrak{D}_{(d),u\,v}^{c\,i},
  \qquad
  d\in\mathcal{J}_{\mathrm{act}},
\]
then define
\[
  \mathcal{O}_{\mathrm{rawact}}^{\E}
  \left[
    {}_{\psi}\mathfrak{D}_{(d),u\,v}^{c\,i}
  \right]
  =
  \E
  \left[
    \exp(X)-1
  \right].
\]
For the diagonal Gaussian mixture marginal, this is available in closed form:
\[
  \widehat{a}_{e}^{c\,i\,d}
  =
  \sum_{k=1}^{K}
  \alpha_{e}^{c\,i\,k}
  \left[
    \exp
    \left(
      \mu_{e}^{c\,i\,k\,d}
      +
      \frac{1}{2}
      (\sigma_{e}^{c\,i\,k\,d})^2
    \right)
    -
    1
  \right]
  \in\R.
\]
Because a Gaussian mixture over $X$ has support on all real values, an
implementation may protect the raw-activity readout by projecting to the
nonnegative line:
\[
  \widehat{a}_{e,+}^{c\,i\,d}
  =
  \max
  \left(
    0,
    \widehat{a}_{e}^{c\,i\,d}
  \right)
  \in\Rnonneg.
\]
The resulting raw-activity field is then aggregated by the endpoint operators
defined in NodeLift.

\subsection{NodeLift Readout Observer}
\label{subsec:nodelift-readout-observer}

The NodeLift procedure can itself be viewed as a graph-level observer.  It reads
an observed input edge feature field, graph matrices, masks, precision weights,
and a synthetic reference gauge, then returns a node feature field and a
residual sidecar.

For one channel, one position, and one anchor time, write
\[
  \mathcal{O}_{\mathrm{NodeLift}}
  \left[
    \left\{
      \tensor{x}_{e}^{c\,i}(t)
      :
      e\in\mathcal{E}_{\mathrm{use}}^{c\,i}(t)
    \right\},
    \mat{A},
    \mat{B},
    \mat{Q},
    \vect{w}_{\star}
  \right]
  =
	  \left(
	    \tensor{s}^{c\,i}(t),
	    {}_{s}\tensor{m}^{c\,i}(t),
	    \tensor{\eps}_{\mathrm{price}}^{c\,i}(t),
	    {}_{\eps}\tensor{m}_{\mathrm{price}}^{c\,i}(t)
	  \right).
\]
Here
\[
  \tensor{s}^{c\,i}(t)
  \in
  \R^{N\times D_y}
\]
is the node feature matrix,
\[
  {}_{s}\tensor{m}^{c\,i}(t)
  \in
  \{0,1\}^{N\times D_y}
\]
is the coordinatewise node validity mask, and
\[
  \tensor{\eps}_{\mathrm{price}}^{c\,i}(t)
\]
is the residual sidecar for price-like coordinates.  A coarse node mask in
$\{0,1\}^{N}$ may be derived from ${}_{s}\tensor{m}^{c\,i}(t)$ when a downstream
module cannot consume coordinatewise validity.

Over a full observed input window, this observer gives
\[
  \tensor{s}_{x}(t)
  \in
  \R^{C\times H_x\times N\times D_y}.
\]
This is the input-side NodeLift observer used by the node-first path.

\subsection{Residual and Cycle-Tension Observers}
\label{subsec:residual-and-cycle-tension-observers}

For a lifted price coordinate, the residual edge field is
\[
  \vect{\eps}_{\mathrm{price}}^{c\,i\,d}
  =
  \vect{x}_{\mathcal{E}}^{c\,i\,d}
  -
  \mat{A}
  \widehat{\vect{y}}^{c\,i\,d}
  \in
  \R^{L_{\mathrm{use}}},
  \qquad
  d\in\mathcal{J}_{\mathrm{price}}.
\]
The residual observer is
\[
  \mathcal{O}_{\mathrm{res}}
  \left[
    \vect{x}_{\mathcal{E}}^{c\,i\,d},
    \widehat{\vect{y}}^{c\,i\,d}
  \right]
  =
  \vect{\eps}_{\mathrm{price}}^{c\,i\,d}.
\]
Given an edge precision matrix $\mat{W}^{c\,i\,d}$, the scalar tension observer is
\[
  \mathcal{O}_{\mathrm{tension}}
  \left[
    \vect{\eps}_{\mathrm{price}}^{c\,i\,d},
    \mat{W}^{c\,i\,d}
  \right]
  =
  \left(
    \vect{\eps}_{\mathrm{price}}^{c\,i\,d}
  \right)^{\top}
  \mat{W}^{c\,i\,d}
  \vect{\eps}_{\mathrm{price}}^{c\,i\,d}
  \in
  \Rnonneg.
\]

This residual is gauge-invariant.  Changing the synthetic reference changes the
displayed node coordinates, but it does not erase the edge-field disagreement
that cannot be explained by one node-potential field.  Later trading logic may
read this residual as cycle tension or arbitrage-like disagreement, but
execution decisions remain outside this document.

\subsection{Synthetic-Reference Observer}
\label{subsec:synthetic-reference-observer}

The synthetic reference used by NodeLift may also be read by an observer.  Let
\[
  \mat{C}_{y}^{\mathrm{quot}}
  \in
  \R^{N\times N}
\]
be a quotient-space node covariance estimate.  The covariance-minimizing
reference observer is
\[
  \mathcal{O}_{\star}
  \left[
    \mat{C}_{y}^{\mathrm{quot}},
    \mat{D}_{\mathrm{node}},
    \vect{w}_0,
    \lambda_{\mathrm{ref}}
  \right]
  =
  \vect{w}_{\star}
  \in
  \mathcal{W}_{\star}.
\]
This observer chooses the synthetic reference basket used to display recovered
node coordinates.  It does not change the edge fit and does not remove residual
cycle tension.

\subsection{Step and Cumulative Return Readouts}
\label{subsec:step-and-cumulative-return-readouts}

The scalar edge-return marginal gives a step return at each forecast position.
For edge-to-node recovery at one forecast step, this step return may be the
desired object:
\[
  r_e^{c\,i}\in\R.
\]
For horizon-level valuation, one may instead read the cumulative log return from
the anchor to forecast step $i$:
\[
  R_e^{c\,i}
  =
  \sum_{j=1}^{i}
  r_e^{c\,j}
  \in\R.
\]
At the level of expectations, the cumulative readout gives
\[
  \widehat{R}_e^{c\,i}
  =
  \sum_{j=1}^{i}
  \widehat{r}_e^{c\,j}
  \in\R.
\]
This cumulative readout is an observer choice%
\note{
  The current MDN head gives horizon-wise marginal distributions.  Summing
  expected step returns is well defined, but the full uncertainty of
  $R_e^{c\,i}$ requires cross-horizon covariance.  Without such a joint horizon
  model, cumulative-return uncertainty must be treated as an observer
  approximation or supplied by an additional covariance model.
}.

The same idea may be applied after NodeLift.  A cumulative node-return readout is
\[
  \widehat{Y}_{u}^{c\,i}
  =
  \sum_{j=1}^{i}
  \widehat{y}_{u}^{c\,j}
  \in\R.
\]
As with cumulative edge returns, the mean is easy to sum, while the uncertainty
requires a cross-horizon covariance model or a scenario-based observer.

\subsection{Trajectory Readouts}
\label{subsec:trajectory-readouts}

The observers above often read one predictive distribution at a time.  We may
also read several forecast horizons together.  These are still readouts of the
MDN field; they do not change the predictive distribution learned by the model.

For one edge $e=(u,v)$ and channel $c$, define a forecast trajectory as
\[
  \vect{F}_{e}^{c,1:H_f}
  =
  \left(
    \vect{f}_{e}^{c\,1},
    \vect{f}_{e}^{c\,2},
    \ldots,
    \vect{f}_{e}^{c\,H_f}
  \right)
  \in
  \left(\R^{D_f}\right)^{H_f}.
\]
The current MDN head gives one marginal predictive distribution
${}_{\psi}\mathfrak{D}_{u\,v}^{c\,i}$ for each horizon $i$.  Therefore, any
trajectory readout must declare how these horizon-wise distributions are
combined.

A simple first choice is the product score
\[
  S_{\mathrm{path}}
  \left(
    \vect{F}_{e}^{c,1:H_f}
  \right)
  =
  \sum_{i=1}^{H_f}
  \log
  {}_{\psi}p_{u\,v}^{c\,i}
  \left(
    \vect{f}_{e}^{c\,i}
    \given
    \tensor{z}_{u\,v}
  \right),
\]
which treats the horizon-wise MDN outputs as conditionally independent given
the edge representation.  This is a readout assumption, not an additional claim
made by the MDN head.

A maximum-density predictive trajectory readout may be written as
\[
  \widehat{\vect{F}}_{e,\mathrm{dens}}^{c,1:H_f}
  =
  \operatorname*{arg\,max}_{\vect{F}_{e}^{c,1:H_f}}
  \left\{
    S_{\mathrm{path}}
    \left(
      \vect{F}_{e}^{c,1:H_f}
    \right)
    -
    \mathcal{R}_{\mathrm{path}}
    \left(
      \vect{F}_{e}^{c,1:H_f}
    \right)
  \right\}.
\]
Here $\mathcal{R}_{\mathrm{path}}$ is an optional trajectory regularizer.  It
may penalize impossible jumps, prefer smooth paths, or enforce consistency
constraints inherited from the market representation.

For implementation, a lighter component-path version is often enough.  For each
horizon and mixture component, define the component-mode score
\[
  s_{e}^{c\,i\,k}
  =
  \log \alpha_{e}^{c\,i\,k}
  -
  \frac{1}{2}
  \log
  \det
  \left(
    2\pi \mat{\Sigma}_{e}^{c\,i\,k}
  \right).
\]
This score measures the height of component $k$ at its own mean.  A discrete
component trajectory is
\[
  \widehat{\kappa}_{e}^{c}
  =
  \operatorname*{arg\,max}_{(k_1,\ldots,k_{H_f})}
  \left\{
    \sum_{i=1}^{H_f}
    s_{e}^{c\,i\,k_i}
    -
    \lambda_{\Delta}
    \sum_{i=2}^{H_f}
    \left\|
      \vect{\mu}_{e}^{c\,i\,k_i}
      -
      \vect{\mu}_{e}^{c\,(i-1)\,k_{i-1}}
    \right\|_2^2
  \right\}.
\]
The resulting representative trajectory is
\[
  \widehat{\vect{F}}_{e,\mathrm{comp}}^{c,1:H_f}
  =
  \left(
    \vect{\mu}_{e}^{c\,1\,\widehat{k}_1},
    \vect{\mu}_{e}^{c\,2\,\widehat{k}_2},
    \ldots,
    \vect{\mu}_{e}^{c\,H_f\,\widehat{k}_{H_f}}
  \right).
\]
This gives one concrete path through the mixture field without collapsing the
whole predictive distribution into an expectation.

\subsection{Scenario Readouts}
\label{subsec:scenario-readouts}

The same distribution field can produce scenario trajectories.  Under the same
horizon-wise product assumption,
\[
  \vect{F}_{e,\mathrm{sc}}^{c,1:H_f,\rho}
  \sim
  \prod_{i=1}^{H_f}
  {}_{\psi}\mathfrak{D}_{u\,v}^{c\,i},
  \qquad
  \rho=1,\ldots,S.
\]
Scenario trajectories are useful when an observer needs more than a mean path,
for example when estimating tail exposure, downside risk, or cumulative-return
uncertainty.  If a later experiment extends NodeLift beyond observed input
windows, it should name the scenario source field explicitly before applying the
graph lift.

\subsection{Control-Prepared Readouts}
\label{subsec:control-prepared-readouts}

The distribution field may also be prepared for a later control problem.  Such a
readout does not solve the trading problem inside this document.  It only
converts predictive distributions, input-side node-lifted fields, covariance
estimates, and residual diagnostics into quantities that a control layer could
consume.

A schematic control-prepared observer may have the form
\[
  \mathcal{O}_{\mathrm{ctrl}}
  \left[
    {}_{\psi}\mathfrak{D}_{\mathcal{E}}^{c\,i}(t),
    \tensor{s}_{x}^{c\,i}(t),
    \tensor{\eps}_{\mathrm{price}}^{c\,i}(t),
    \vect{h}(t)
  \right]
  =
  \left(
    \tensor{s}_{x}^{c\,i}(t),
    \mat{C}_{y}^{c\,i}(t),
    \text{risk fields},
    \text{residual fields},
    \text{constraint fields}
  \right).
\]
These objects are compatible with later dynamic-programming or
Hamilton--Jacobi--Bellman formulations, but they are not the policy itself.

In that later setting, one would introduce a control state $\mathfrak{s}$, a
control $a$, a value function $V$, and an objective.  A schematic
continuous-time control equation would have the form
\[
  0
  =
  \partial_t V(\mathfrak{s},t)
  +
  \sup_{a\in\mathcal{A}(\mathfrak{s})}
  \left\{
    \operatorname{Reward}(\mathfrak{s},a)
    -
    \mathrm{Cost}(\mathfrak{s},a)
    +
    {}_{\psi}\mathcal{L}^{a}V(\mathfrak{s},t)
  \right\}.
\]
The operator ${}_{\psi}\mathcal{L}^{a}$ would be built from drift, covariance,
risk, node-lifted state, and residual-tension quantities read from the
forecasting interface.  The control problem remains a separate layer.

\subsection{Second-Order Observers and Normalization}
\label{subsec:second-order-observers-and-normalization}

If a joint second-order observer is maintained across a collection of edge,
node, residual, or trajectory readouts
\[
  Z_1,\ldots,Z_m\in\R,
\]
define the covariance matrix by
\[
  \Sigma_{a\,b}
  =
  \Cov(Z_a,Z_b),
  \qquad
  \Sigma\in\R^{m\times m}.
\]
The diagonal terms $\Sigma_{a\,a}$ are variances.  The normalized correlation
matrix is
\[
  \rho_{a\,b}
  =
  \frac{
    \Sigma_{a\,b}
  }{
    \sqrt{\Sigma_{a\,a}\Sigma_{b\,b}}
  },
  \qquad
  \rho\in\R^{m\times m}.
\]
This normalization removes scale and keeps only the signed strength of linear
co-movement.

A covariance observer is additional information beyond the pointwise MDN head.
It may come from repeated evaluations, sampled scenarios, an ensemble, a
separate covariance model, NodeLift covariance propagation, or a policy-level
estimator.  The current MDN head standardizes the edge-local predictive
distributions; it does not by itself impose a joint covariance model across the
whole market graph, across channels, or across forecast horizons.

\calloutbox{End of the Observer Layer.}{\\%
  Observers define how uncertainty-bearing edge distributions become usable
  fields: scalar returns, activity readouts, node-lifted tensors, residual
  tension, trajectories, scenarios, covariance summaries, and control-prepared
  quantities.  They read the forecasting interface; they do not replace it.
}
